\documentclass[12pt]{article}
\usepackage[utf8]{inputenc}
\usepackage[T1]{fontenc}
\usepackage[margin=1in]{geometry}
\usepackage{setspace}
\usepackage{xcolor}
\usepackage[colorlinks=true, urlcolor=blue!60!black, linkcolor=blue!60!black]{hyperref}
\setlength{\parindent}{0em}
\setlength{\parskip}{0.45em}
\renewcommand{\baselinestretch}{1.0}
\begin{document}
\begin{center}
{\Large \textit{San Antonio Independent School District v.\ Rodriguez}, 411 U.S.\ 1 (1973)\par}
\vspace{0.3em}
{\normalsize Case Brief --- One-Page Edition\par}
\end{center}

\textbf{Memory Jogger.} Property-tax school funding disparities do not violate equal protection (\href{https://supreme.justia.com/cases/federal/us/411/1/}{Justia}).

\textbf{Facts.} Texas paired a state program guaranteeing basic education with capped local property taxes. Property-poor, largely Mexican-American Edgewood spent \$356 per pupil; affluent Alamo Heights spent \$594. Edgewood parents sued.

\textbf{Procedural History.} A three-judge court struck down the system; Texas appealed directly.

\textbf{Votes.} Reversed, 5--4: Powell for the Court; Stewart concurred; Brennan, White, and Marshall dissented.

\textbf{Holding.} District wealth is not a suspect class and education is not a fundamental right, so rational basis review applies, and local control satisfies it.

\textbf{Reasoning.} No definable class of indigents suffered an absolute deprivation; a right is fundamental only if the Constitution guarantees it, and education's importance does not suffice. White found that capped rates gave poor districts no real choice; Marshall urged a ``spectrum of standards.''

\textbf{Examples.} \emph{Same-Side:} school-bus fees (\textit{Kadrmas}). \emph{Opposite-Side:} excluding undocumented children (\textit{Plyler}). \emph{Fence-Sitter:} a claimed right to basic literacy.

\textbf{Critique.} Defenders credit state-court reform; progressives say it entrenched unequal schools.

\textbf{Key Quotation.} ``The system of alleged discrimination and the class it defines have none of the traditional indicia of suspectness: the class is not saddled with such disabilities, or subjected to such a history of purposeful unequal treatment, or relegated to such a position of political powerlessness as to command extraordinary protection from the majoritarian political process.'' --- Justice Powell, for the Court.

\end{document}
